\documentclass[11pt]{article}

\usepackage[T1]{fontenc}
\usepackage{lmodern}
\usepackage[margin=1in]{geometry}
\usepackage{amsmath}
\usepackage{xcolor}
\usepackage{listings}
\usepackage{cleveref}

% Custom colors for the source-listing style.
\definecolor{codekeyword}{rgb}{0.10,0.10,0.65}
\definecolor{codecomment}{rgb}{0.30,0.50,0.30}
\definecolor{codestring}{rgb}{0.62,0.13,0.20}
\definecolor{codeframe}{gray}{0.55}
\definecolor{codeback}{gray}{0.97}

% A named listing style: colored keywords/comments/strings, a single frame,
% left-hand line numbers, and the caption printed below the float.
\lstdefinestyle{pystyle}{
  language=Python,
  backgroundcolor=\color{codeback},
  basicstyle=\ttfamily\footnotesize,
  keywordstyle=\bfseries\color{codekeyword},
  commentstyle=\itshape\color{codecomment},
  stringstyle=\color{codestring},
  numbers=left,
  numberstyle=\tiny\color{codeframe},
  numbersep=8pt,
  stepnumber=1,
  frame=single,
  rulecolor=\color{codeframe},
  framesep=4pt,
  captionpos=b,
  showstringspaces=false,
  breaklines=true,
  columns=fullflexible,
  tabsize=4,
  aboveskip=1.2em,
  belowskip=1.2em
}
\lstset{style=pystyle}

% Make cleveref name the lstlisting float a "Listing".
\crefname{lstlisting}{Listing}{Listings}
\Crefname{lstlisting}{Listing}{Listings}

\title{Reproducible Source Listings as a Conversion Stress Case}
\author{Test Corpus}
\date{}

\begin{document}
\maketitle

\begin{abstract}
This fixture exercises the \texttt{listings} package under conditions that
commonly appear in a reproducibility appendix: a custom, colored listing style,
an inline code span, a displayed listing typed directly into the source, and a
second listing read from a child file on disk. The accompanying prose carries
cross-references to the code so that a faithful conversion must preserve both
the formatted blocks and the numbered references that point at them.
\end{abstract}

\section{Motivation}
\label{sec:motivation}

Empirical papers increasingly ship the code that produced their estimates, and
that code is usually presented through a package such as \texttt{listings} rather
than as ordinary prose. The distinguishing features are syntactic: keywords,
comments, and string literals are colored differently, leading whitespace is
preserved, and each physical line carries a number so that the surrounding text
can refer to it. A converter that flattens these elements into a single
paragraph loses the indentation that makes the code legible and discards the line
numbering on which the discussion depends. The remainder of this note assembles
the constructs that most often expose such failures.

Inline code appears in running text as well. When a function such as
\lstinline|standardized_car(abnormal, forecast_sd)| is named mid-sentence, it
should remain monospaced and visually distinct from the body font, exactly as an
identifier like \lstinline|forecast_sd| should. Inline spans share the typeface
of the displayed blocks but carry neither a frame nor a line number.

\section{A Typed Listing}
\label{sec:typed}

\Cref{lst:accumulate} is entered directly into the document source. It contains a
docstring, an inline comment, a string literal, and four levels of indentation,
all of which the listing style renders with distinct treatment. The line numbers
on the left margin let the discussion single out individual statements.

\begin{lstlisting}[caption={Accumulating an abnormal return over the event window},label={lst:accumulate}]
def accumulate(abnormal_returns):
    """Sum abnormal returns across the event window."""
    total = 0.0
    for value in abnormal_returns:
        if value is None:          # skip non-trading days
            continue
        total += value
    return total
\end{lstlisting}

The loop in \cref{lst:accumulate} guards against missing observations before it
adds each abnormal return to the running total, a pattern familiar from
event-study code that must tolerate non-trading days. The standardization that
typically follows this accumulation appears separately in \cref{lst:external}.

\section{A Listing Read From Disk}
\label{sec:external}

The second listing is not present in this file. It is pulled in at compile time
with \verb|\lstinputlisting| from \texttt{code/sample.py}, so the converter must
resolve the external path and apply the same named style that governs the typed
block above. \Cref{lst:external} reports the standardization step that scales the
cumulative abnormal return by its forecast standard deviation.

\lstinputlisting[style=pystyle,caption={Standardizing the cumulative abnormal return (read from \texttt{code/sample.py})},label={lst:external}]{code/sample.py}

Read together, \cref{lst:accumulate,lst:external} trace a small but complete
computation: the first accumulates the raw signal and the second rescales it into
a test statistic. Because one listing is typed inline and the other is read from
a child file, a faithful conversion must handle both ingestion paths while
keeping the colored syntax, the line numbers, the captions printed below each
block, and the cross-references that bind the prose to the code.

\section{Discussion}
\label{sec:discussion}

The constructs gathered here are deliberately ordinary; each appears routinely in
computational appendices. Their combination is what makes the fixture demanding,
since a converter must simultaneously honor a user-defined style, an inline code
span, a directly typed float, an externally included float, and cleveref
references that name the floats as listings rather than as generic figures.
\Cref{sec:motivation,sec:typed,sec:external} together establish that the code is
embedded in an otherwise conventional article with sectioning and prose.

\end{document}
